When the particles formed in the reaction are inside the ``reaction region'', the reaction probability is proportional to the squared modulus of their wave function. This region corresponds to separations of the order of the nuclear-force range $a$ (compare the analogous argument in \S~143 for the initial particles). Here we require only the dependence of the reaction probability on the characteristics of the relative motion of one nucleon pair. It is therefore sufficient to consider the wave function $\psi_{\mathbf p}(\mathbf r)$ of this motion alone. Thus, the probability of producing a nucleon pair with relative momentum in the interval $d^3p$ is
